\section{From Fixed Cost to Recurring Tax}
\label{sec:tax}

Classic economic geography explains why distance matters even when transport
is cheap. In \citeauthor{krugman1991}'s \parencite*{krugman1991} model, firms
cluster near large markets to exploit economies of scale, and once a core
forms it attracts more activity: the periphery is not just far away but
\emph{left behind} by a self-reinforcing process. \textcite{limao2001} show
how much geography and infrastructure matter: poor infrastructure and
landlocked status can multiply transport costs, and trade volumes fall
steeply as costs rise. \textcite{briguglio1995} catalogued the specific
vulnerabilities of small islands: high per-unit costs of infrastructure and
administration, dependence on a narrow range of exports, exposure to shocks
they cannot influence. And the expected \enquote{death of distance} has been
slow to show up in the data: the estimated effect of distance on trade has
not declined \parencite{disdier2008}, even as transport technology improved
\parencite{hummels2007}.

What changes in an integrated, accelerating world is the \emph{frequency}
with which these costs are paid. A remote economy in 1960 paid a high freight
rate on a few bulk exports and imports. A remote economy in 2025 pays for
every parcel, every software subscription, every cross-border payment, every
specialist consultation and every round of trade or climate negotiations.
The cost is no longer an entry fee; it is a subscription.

\begin{keyidea}[The recurring tax]
Remoteness becomes a \emph{recurring tax on participation} when economic life
consists of many small, frequent, time-sensitive transactions rather than a
few large ones. The tax is levied in freight and insurance, in bandwidth and
latency, in the scarcity of specialists, and in the cost of being present
where rules are made.
\end{keyidea}

\Cref{fig:engine} sketches why the taxes \emph{compound}. Higher trade costs
raise the cost of inputs and shrink the market, weakening the return on
digital investment. Thin and fragile connectivity makes it easier for
skilled people to work for distant employers---or to leave. Their departure
thins the institutions that should negotiate better freight contracts,
spectrum deals, fishing licences and climate finance. Each link is modest; the
loop is what turns geography into a structural gap.

% The "compounding engine": four recurring taxes that feed one another.
\begin{figure}[t]
\centering
\begin{tikzpicture}[
    node distance=1.2cm and 1.6cm,
    tax/.style={rectangle, rounded corners=3pt, draw=teal, line width=0.9pt,
      fill=mist, align=center, font=\sffamily\small, text width=3.3cm,
      minimum height=1.25cm, inner sep=5pt},
    hub/.style={circle, fill=ember, text=white, align=center,
      font=\sffamily\bfseries\footnotesize, minimum size=2.35cm, inner sep=1pt},
    loop/.style={-{Stealth[length=3mm]}, line width=1.2pt, teal, bend left=18},
    lab/.style={font=\sffamily\scriptsize\itshape, text=slate, align=center, fill=white,
      inner sep=1.5pt},
  ]
  \node[hub] (core) {Small, far\\ \& thinly\\ connected};
  \node[tax, above=1.05cm of core] (trade) {\textbf{Trade tax}\\ fewer ships, higher freight, thinner markets};
  \node[tax, right=1.25cm of core] (tech) {\textbf{Connectivity tax}\\ one cable, costly bandwidth, slow repair};
  \node[tax, below=1.05cm of core] (people) {\textbf{Talent tax}\\ skilled emigration, thin specialist cover};
  \node[tax, left=1.25cm of core] (inst) {\textbf{Institutional tax}\\ small bureaucracies, weak bargaining power};

  \draw[loop] (trade.east) to node[lab, above right=-2pt] {costlier inputs,\\ weaker returns to\\ digital investment} (tech.north);
  \draw[loop] (tech.south) to node[lab, below right=-2pt] {remote work pulls\\ talent toward\\ the centre} (people.east);
  \draw[loop] (people.west) to node[lab, below left=-2pt] {fewer specialists\\ to run the state} (inst.south);
  \draw[loop] (inst.north) to node[lab, above left=-2pt] {weaker leverage in\\ trade and finance\\ negotiations} (trade.west);

  \foreach \n in {trade,tech,people,inst}
    \draw[-{Stealth[length=2mm]}, ember, line width=0.7pt, dashed] (core) -- (\n);
\end{tikzpicture}
\caption{\takeaway{The taxes compound.} Each recurring cost of remoteness
  raises the others: the loop, not any single link, is what turns a
  geographic fact into a structural gap. Dashed arrows: smallness and distance
  feed every channel directly.
  \source{Author's schematic, drawing on \textcite{krugman1991,limao2001,briguglio1995}.}}
\label{fig:engine}
\end{figure}


The loop is a hypothesis, not a finding, and the sections that follow test
each link with data. Two things should be held in mind. First, every arrow can
run in reverse: remittances flow back, diasporas carry knowledge home, and
remote places can earn rents from their position (\cref{sec:counter}).
Second, the loop does not say remote places get poorer in absolute terms. It
says their position becomes harder to improve, more volatile and less within
their own control.
